% Catálogo de componentes del taller (taller-recursos.sty): un ejemplo de cada uno.
% Compilar: bash curriculo/plantilla/compilar.sh curriculo/plantilla/ejemplos/catalogo-taller.tex
\documentclass[taller]{planeacion}
\graphicspath{{../}}
\materia{filosofia}
\grado{10.°} \periodo{I} \anio{2027}
\begin{document}
\fichataller{Catálogo}{Componentes del taller}{55 min}{En parejas}
\titulotaller{Componentes del taller}{Un ejemplo de cada comando; la versión del docente muestra las soluciones en vino.}
\seccion{1}{Textos, fuentes y recursos}
\texto{TEXTO 1 · RÓTULO}{Contenido del texto citado, con el fragmento exacto.}{Autor, \emph{Obra} (año), cap. o p. (traducción propia)}
\recuerda{Explicación breve con un ejemplo, para que el taller no dependa de la memoria.}
\concepto{Proposición}{Enunciado que puede ser verdadero o falso.}
\recurso{VIDEO}{Título del recurso (autor, año, 4 min)}{Qué hacer mientras lo ven. Sin enlace sale «QR pendiente».}{}
\figura[0.3]{no-existe.png}{Pie con autor, obra, año y licencia. (Archivo faltante: recuadro de aviso.)}
\imagenlado[0.25]{no-existe.png}{Imagen a la izquierda y su explicación a la derecha (\texttt{\textbackslash imagenlado}).}
\noindent\euler{\draw (0,0) circle (1.2); \draw[fill=black!12] (0,-0.25) circle (0.55); \node at (0,-0.25) {S}; \node at (0,0.8) {P};}{Diagrama de Euler (\texttt{\textbackslash euler})}\par
\respuestaia{Respuesta de un asistente de IA con un error plausible para analizar.}
\seccion{2}{Actividades}
\reto{0}{Punto con tiempo y puntaje}{10 min}{2,0}
Completar en línea: la proposición «Algún S no es P» es de tipo \hueco[O]{1.5cm}.
\begin{verdaderofalso}
  \vf{V}{Toda proposición A es universal.}
  \vf{F}{La O se convierte simplemente.}
\end{verdaderofalso}
\begin{opciones}
  \op{Conversión de una A}
  \op*{Obversión de una E}
\end{opciones}
\pista{Ayuda breve que orienta sin regalar la respuesta.}
\reto{0}{Relacionar}{10 min}{1,5}
\relacionar{{Contrarias},{Contradictorias},{Subalternas}}{{A y O},{A e I},{A y E}}[1/c,2/a,3/b]
\reto{0}{Tabla y respuesta abierta}{10 min}{1,5}
\begin{tablataller}{|>{\RaggedRight\arraybackslash}p{5cm}|>{\RaggedRight\arraybackslash}X|}
  \cabecera\tit{Dato} & \tit{Respuesta}\\ \hline
  Ejemplo & \sol{Solución del docente}\\ \hline
\end{tablataller}
\respuesta[2]{Respuesta modelo que solo ve el docente.}
\cierre{5,0}{\logro{Logro 1}\logro{Logro 2}}
\end{document}
